\begin{table}[htbp]
\centering
\caption{Diagnóstico das covariáveis imputadas}
\label{tab:diagnostico_imputacao}
\begin{tabular}{lcccc}
\toprule
Série-disciplina & N completas & N imputadas & Baseline comp. & Baseline imput. \\
\midrule
9EF LP & 13.992 & 15.487 & 230,07 & 229,75 \\
9EF Mat & 13.999 & 15.495 & 246,56 & 245,87 \\
EM3 LP & 1.032 & 27.842 & 266,68 & 265,49 \\
EM3 Mat & 1.033 & 27.857 & 272,48 & 269,72 \\
\bottomrule
\end{tabular}
\vspace{0.2cm}
\begin{flushleft}
\footnotesize Nota: observações completas possuem todas as covariáveis brutas usadas no DR-DiD sem imputação. Diferenças de baseline entre colunas indicam que o diagnóstico sem imputação também muda a composição da amostra. Baseline é a proficiência média na escala SARESP.
\end{flushleft}
\end{table}
